\chapter{Simulation configurations and batch workflows}\label{ch:workflows}
\section{Three different meanings of configuration}
An OpenRocket design can carry multiple motor and staging arrangements, multiple experiments on those arrangements, and multiple random realizations of an experiment. Keeping these levels separate is essential when extending the application.

\begin{tabularx}{\linewidth}{@{}L{3.5cm}YY@{}}\toprule
Object & What it selects & What it does not own\\\midrule
\class{FlightConfiguration} & A configuration ID, active stages, installed motors, and configuration-dependent component parameters & Weather, integrator settings, or recorded trajectories\\
\class{Simulation} and \class{SimulationOptions} & One selected flight configuration plus launch conditions, models, timing, extensions, and ensemble settings & A separate editable copy of the entire design for every saved simulation\\
\class{SimulationConditions} & Runtime calculators, listeners, configuration, and environment prepared for one execution & The GUI's editable settings model\\
\class{EnsembleRunParameters} & The particular environmental and coefficient draws, motor seed, and wind bearing for one member & A second independent ensemble nested inside the member\\\bottomrule
\end{tabularx}

The saved simulation is the experiment's durable handle. It links the design and settings to its latest result. A run translates editable options into runtime conditions and creates mutable execution state. The numerical engine described in Chapter~\ref{ch:engine} consumes this prepared state; it does not reach into Swing controls to obtain values.

Changing design geometry, flight configuration, relevant options, or extension configuration can make a result stale. The fork also tracks flight-computer design provenance and fingerprints. A historical result remains evidence of the settings that produced it, even when the current editor shows different settings. Imported external data has its own status and must not be mistaken for a newly executed simulation.

\sourceline{\source{core/src/main/java/info/openrocket/core/document/Simulation.java}{Simulation}; \source{core/src/main/java/info/openrocket/core/simulation/SimulationOptions.java}{SimulationOptions}; \source{core/src/main/java/info/openrocket/core/simulation/SimulationConditions.java}{SimulationConditions}.}

\section{Configuring an ordinary flight}
Before any batch or control feature is involved, a simulation specifies a vehicle configuration and a launch experiment. The flight-configuration ID selects the motor installations and configuration-specific ignition, separation, and deployment settings in the component tree. The simulation options then provide the launch site, guide, environment, and numerical controls. Two simulations can select the same motors and differ only in weather; two configurations can install different motors while using otherwise identical launch conditions.

\begin{tabularx}{\linewidth}{@{}L{3.1cm}YY@{}}\toprule
Setting family & What it controls & Consequence in the execution path\\\midrule
Launch site & Latitude, longitude, altitude and geodetic method & Initial world position, gravity, atmospheric altitude, coordinate propagation, and optional Coriolis terms\\
Launch guide & Length, inclination and bearing & Initial attitude and constrained travel before rod clearance\\
Atmosphere & Standard atmosphere or launch temperature/pressure & Density, speed of sound and viscosity evaluated along the flight\\
Wind & Average/turbulent or altitude-dependent wind settings and seed & Air-relative velocity, angle of attack and drift; the wind is not simply added to ground speed after integration\\
Numerics & Maximum timestep, angular-step limits and maximum simulation duration & Resolution, event handling and run termination; a smaller maximum step is not itself a new physical model\\
Extensions & Selected features and their saved settings & Runtime listener initialization and extra data channels\\\bottomrule
\end{tabularx}

Defaults for a newly created simulation come from application preferences. Once a scenario has its own options, those are the inputs to its next run. The status displayed beside a simulation distinguishes an unrun case, usable stored data, stale data, imported data, and execution problems. Reusing a result is therefore a decision about the saved definition and its provenance, not merely whether an array of samples exists.

Motor selection also includes ejection delay and ignition policy. A curve's manufacturer/designation describes the motor record, while the installation associates that record with a mount and flight configuration. A clustered mount can produce several motor instances from one configuration entry. The motor chooser, configuration tables, and simulation options edit these different owners rather than merging them into one global motor selection.

\sourceline{\source{core/src/main/java/info/openrocket/core/simulation/SimulationOptions.java}{ordinary simulation options}; \source{core/src/main/java/info/openrocket/core/rocketcomponent/FlightConfiguration.java}{flight configuration}; \source{core/src/main/java/info/openrocket/core/motor/MotorConfiguration.java}{motor installation}.}

\section{Static component analysis and parameter sweeps}\label{sec:static-analysis}
The inherited component-analysis tools explain why the current design has a particular mass, CP, drag coefficient, or roll contribution. They ask the same mass and aerodynamic calculators used elsewhere for a decomposition by component. This is useful before running a trajectory: a large pressure contribution can be traced to a shape or exposed junction, while an unexpected CG can be traced to a material, motor, instance count, or override.

\class{CAParameters} holds the analysis conditions. \class{CAParameterSweep} copies those conditions, varies Mach number, angle of attack, lateral wind/flow angle, or roll rate, constructs \class{FlightConditions}, and asks the aerodynamic calculator for force analysis at every point. \class{CADataBranch} stores the domain values and per-component channels; \class{CADataType} and \class{CADomainDataType} describe the quantities and units. The Swing analysis dialog supplies component selection, tables, plots, and export.

These are controlled static evaluations. Sweeping Mach at fixed angle of attack does not integrate a trajectory through increasing speed, burn propellant, or fly the vehicle through a wind profile. The sweep constructs flight conditions from the selected rocket configuration and analysis parameters; it does not automatically reproduce every atmospheric setting of an arbitrary saved simulation. A comparison with a flight-history sample must align the relevant conditions explicitly.

CP and normal-force slope are related: the combined CP is weighted by component normal-force contributions. A negligible total weighting makes a CP undefined, rather than establishing a meaningful CP at the origin. Static stability is conventionally the axial separation $x_{\rm CP}-x_{\rm CG}$; dividing by reference diameter expresses it in calibers. Motor loading, reference diameter, flow angle, and Mach can change the result. A positive static margin is not a complete prediction of dynamic stability throughout a flight.

Drag analysis distinguishes pressure, base, and friction terms and the distinction between a component's per-instance value and the installed total. A single tree entry can represent several fins, lugs, or repeated assemblies. The mass-analysis map similarly retains contributions and override information rather than treating every row as an unrelated lump.

\begin{contract}
In this snapshot, \class{CAParameterSweep} forms its channel named \code{TOTAL_ROLL_COEFFICIENT} by adding the forcing and damping values. The flight engine's net roll term subtracts damping from forcing. Therefore that analysis channel must not be identified with the engine's net torque coefficient without checking the sign convention. The separate forcing and damping channels are the unambiguous comparison points.
\end{contract}

\sourceline{\source{core/src/main/java/info/openrocket/core/componentanalysis/CAParameterSweep.java}{CAParameterSweep}; \source{core/src/main/java/info/openrocket/core/componentanalysis/CADataBranch.java}{CADataBranch}; \source{swing/src/main/java/info/openrocket/swing/gui/dialogs/componentanalysis/ComponentAnalysisDialog.java}{component-analysis dialog}; \source{core/src/main/java/info/openrocket/core/masscalc/CMAnalysisEntry.java}{mass-analysis entries}.}


\section{Ensemble execution and ownership}
An ensemble repeats an experiment with controlled variation. \class{EnsembleSettings} contains the count, seed, worker limit, and distributions; \class{EnsembleRunner} coordinates members; \class{EnsembleRunListener} applies the member's runtime perturbations; \class{EnsembleAccumulator} produces aggregate statistics; and \class{EnsembleRunArchive} retains compressed individual flights.

\diagram{fig:ensemble}{An ensemble isolates member execution, consumes members in a deterministic order, and publishes a complete result.}{
\begin{tikzpicture}
\simple{parent}{0,0}{Parent simulation\\settings + frozen FC design}{4cm}
\simple{draw}{6,0}{Coordinator\\ordered parameter draws}{4cm}
\simple{members}{6,-2.2}{Bounded worker pool\\isolated rocket and simulation}{4cm}
\simple{acc}{0,-2.2}{Ordered consumption\\streaming statistics}{4cm}
\simple{archive}{0,-4.4}{Compressed archive\\complete member histories}{4cm}
\simple{publish}{6,-4.4}{Final FlightData\\mean + deviations + metrics}{4cm}
\draw[use] (parent)--(draw); \draw[use] (draw)--(members);
\draw[use] (members)--(acc); \draw[use] (acc)--(archive);
\draw[use] (acc)--(publish); \draw[use] (archive)--(publish);
\end{tikzpicture}}

Each member receives a copied rocket with component IDs preserved, a separate document and simulation, its chosen flight configuration, its own random seed, and cloned listeners. The child has ensemble execution disabled. Preserving IDs lets recovery and actuator assignments resolve against the copied geometry; preserving IDs does not imply sharing the component instances.

The coordinator makes random draws in member order before handing work to the pool. Wind/atmosphere/motor choices, coefficient perturbations, and wind-bearing perturbations use the intended random streams so adding a new category need not silently scramble unrelated draws. Results are consumed in member order even if workers finish out of order. This makes aggregate arithmetic reproducible across worker counts and bounds the number of live, uncompressed histories by the worker window. Compressed archives still grow with the amount of retained data.

The default worker count is four; validated settings can select another count, capped at 256 and reduced to the number of runs. Parallel execution rejects simulation extensions other than the specifically supported flight-computer extension. This is a concrete compatibility rule: cloning an arbitrary extension is not proof that its static state or external resources are thread safe. The selected external FC design is frozen for the ensemble so a file edit midway through the run cannot change later members.

Cancellation interrupts workers and closes the executor and archive-building work. The parent publishes an ensemble only after the requested members have produced a valid complete result. A failed, aborted, or structurally inconsistent member is not silently omitted from the denominator.

\section{What an ensemble varies}
Member parameters include launch temperature and pressure, east/north wind components, wind bearing, a motor-noise seed, and multiplicative mass, drag, roll-force, and roll-damping factors. A coefficient draw has the form $m=1+\sigma Z$, with $Z$ a standard normal variate. Factors that must be positive are validated; invalid realizations are rejected rather than clipped into an artificial point mass at zero. Disabled variations reduce to the nominal values and unit multipliers.

The mass and aerodynamic perturbations are applied through simulation callbacks. They express uncertainty in the evaluated model; they are not a permanent edit to the user's component tree. A coefficient multiplier is also a different hypothesis from changing a geometric dimension and recomputing all dependent properties. The settings report should preserve which hypothesis was selected.

\class{GaussianThrustNoise} is evaluated from seeded knots indexed by physical time and motor identity. It is not a new random number on every Runge--Kutta derivative call. Adjacent normal samples are interpolated and divided by $\sqrt{(1-f)^2+f^2}$, where $f$ is the interpolation fraction, preserving marginal variance during interpolation. The resulting process is time-correlated, with its behavior set by the knot spacing. During the burn, noisy thrust is clamped to a nonnegative value. During thrusting, the ensemble limiter caps a step at one quarter of the configured noise interval. Internal RK evaluations and worker scheduling do not advance a mutable random stream differently.

\sourceline{\source{core/src/main/java/info/openrocket/core/simulation/ensemble/EnsembleSettings.java}{EnsembleSettings}; \source{core/src/main/java/info/openrocket/core/simulation/ensemble/EnsembleRunner.java}{EnsembleRunner}; \source{core/src/main/java/info/openrocket/core/simulation/ensemble/EnsembleRunListener.java}{EnsembleRunListener}; \source{core/src/main/java/info/openrocket/core/simulation/ensemble/GaussianThrustNoise.java}{GaussianThrustNoise}.}

\section{Statistics, alignment, and interpretation}
Member timesteps need not coincide. The accumulator starts with the first member's time grid and restricts output to the physical-time interval common to all members, including their ground-hit limits. It aligns subsequent channels onto that grid and tracks streaming moments. Continuous scalar channels use their corresponding interpolated values; discrete and special channels require their own treatment. A missing value is not substituted with zero.

Ordinary scalar bands describe the sample standard deviation, with the $n-1$ denominator, around the ensemble mean. Wrapped angles use circular treatment. Quaternion samples are aligned to a common hemisphere to avoid cancellation between $q$ and $-q$; the playback reader normalizes the component mean before using it as an orientation. A mixed FC state is marked mixed, not presented as an intermediate state-machine state. Recovery channels can express the fraction of members deployed at a particular physical time. These summaries answer different questions and should not be plotted as if they all had ordinary continuous semantics.

Only matching events present in all members enter the mean event series, at their average time, and only inside the common output interval. Scalar metrics such as each member's apogee are collected from individual full flights. Thus the average of member apogees is not generally the peak of the mean altitude curve. Individual archived histories remain available beyond the common interval, so a shorter member does not destroy the longer member's descent record.

\begin{contract}
An ensemble is a model-based sensitivity experiment. Its bands reflect the selected distributions and shared model assumptions. They are not automatically calibrated confidence intervals for a real vehicle, and a mean trajectory is not a trajectory flown by a particular member.
\end{contract}
\sourceline{\source{core/src/main/java/info/openrocket/core/simulation/ensemble/EnsembleAccumulator.java}{EnsembleAccumulator}; \source{core/src/main/java/info/openrocket/core/simulation/ensemble/EnsembleResult.java}{EnsembleResult}; \source{core/src/main/java/info/openrocket/core/simulation/ensemble/EnsembleRunArchive.java}{EnsembleRunArchive}.}

\section{Motor comparison batches}
Motor comparison asks how a fixed design behaves with different candidate motors. It is a separate workflow from random variation of one motor. The GUI gathers a baseline simulation, motor mount, candidate filters, constraints, and delay policy into a snapshot. \class{MotorCandidateSelector} filters the motor database; \class{MotorClearance} checks fit and intersections; \class{MotorBatchRunner} executes isolated cases; metric and report classes retain the outcome of each case.

Filters cover dimensions and motor characteristics such as manufacturer, impulse, thrust, burn duration, and certification information. Diameter fit and geometric intersections are distinct checks. An overlength motor can be represented as an advisory case rather than categorically excluded; an intersecting installation requires the workflow's explicit override and retains its warning. This distinction allows intentional motor protrusion without hiding collisions with other components.

The runner uses a bounded pool, copies the baseline for individual cases, applies the chosen motor and delay, and records failures alongside successful results. The listed-delay search evaluates candidate delays and ranks eligible outcomes using deployment metrics. A batch report is not the same object as an ensemble mean: it is a comparison table whose rows keep their motor identities, warnings, constraints, and outcomes. Selected results can become ordinary simulations for further inspection. CSV exports and saved report history support reviewing the comparison after the worker has finished.

\sourceline{\source{core/src/main/java/info/openrocket/core/simulation/motorbatch/MotorBatchSnapshot.java}{MotorBatchSnapshot}; \source{core/src/main/java/info/openrocket/core/simulation/motorbatch/MotorCandidateSelector.java}{MotorCandidateSelector}; \source{core/src/main/java/info/openrocket/core/simulation/motorbatch/MotorClearance.java}{MotorClearance}; \source{core/src/main/java/info/openrocket/core/simulation/motorbatch/MotorBatchRunner.java}{MotorBatchRunner}; \source{swing/src/main/java/info/openrocket/swing/gui/simulation/MotorBatchDialog.java}{MotorBatchDialog}.}

\section{Design optimization}
The inherited optimization framework changes selected design variables to improve an objective. \class{SimulationModifier} implementations define the variables and their ranges. An \class{OptimizableParameter} extracts a quantity from the trial design or simulated result, while an \class{OptimizationGoal} converts it into a minimization objective. Domain constraints penalize unacceptable trial designs. The numerical search layer supplies algorithms independently of the rocket-specific objective adapter.

\class{RocketOptimizationFunction} duplicates the rocket and simulation for an evaluation, applies the trial modifiers, checks the domain, and evaluates the requested quantity. The GUI's optimization worker coordinates progress and plots. A one-variable search and a multidimensional search can use different algorithms without teaching the flight engine how optimization works.

These workflows can share design and result types, but their semantics differ: an ordinary simulation predicts one configured flight; an ensemble summarizes randomized members; a motor batch compares named alternatives; an optimizer searches a variable space. New UI or export code should retain that distinction rather than flattening all four into an undifferentiated list of trajectories.
\sourceline{\source{core/src/main/java/info/openrocket/core/optimization/rocketoptimization/RocketOptimizationFunction.java}{RocketOptimizationFunction}; \source{core/src/main/java/info/openrocket/core/optimization/services/OptimizationServiceHelper.java}{OptimizationServiceHelper}; \source{swing/src/main/java/info/openrocket/swing/gui/dialogs/optimization/GeneralOptimizationDialog.java}{GeneralOptimizationDialog}.}

\section{Optimization variables, goals, and search mechanics}
The optimization services offer a bounded set of meaningful design variables. The default modifier service includes nose and transition dimensions, body-tube dimensions, fin chords/span/sweep/cant where supported, launch-lug dimensions, explicit component mass, parachute diameter and drag coefficient, and streamer dimensions/aspect ratio/drag coefficient. Availability depends on component type, automatic sizing, fill state, and whether a coefficient is automatic. It avoids offering incompatible parameterizations indiscriminately, such as two independent controls for quantities already linked by geometry.

\class{GenericComponentModifier} targets a component by identity and applies a property change through the appropriate model API. Configuration modifiers handle values associated with a flight configuration. The search variables have physical bounds and display units, but the numerical optimizer operates in its own normalized variable space. A variable being offered in the dialog does not mean the search can violate the underlying component's geometry constraints.

Objective services include maximum altitude, speed and acceleration; absolute or relative stability; impact velocity; landing distance; flight time; and deployment velocity. Some quantities can be obtained from static calculations, while trajectory objectives require a simulation. A minimization goal preserves the objective direction, a maximization goal reverses it for the minimizer, and a value-seeking goal rewards proximity to a target. The stability domain can restrict acceptable designs; the identity domain supplies no such additional restriction.

The general optimizer layer contains one-dimensional golden-section search and multidirectional search for several variables, along with points, controllers, caches, and stopping/progress interfaces. Those classes know about functions and evaluations, not nose cones or launch rods. \class{RocketOptimizationFunction} is the adapter that copies the design, applies modifiers, checks the domain, and evaluates the requested physical quantity. The GUI worker presents progress and candidate designs without running the search on Swing's event-dispatch thread.

The returned design is the search result for the selected variables, bounds, goal, and model. It does not imply a global optimum over all possible rocket shapes, especially when the feasible region is constrained or the objective is irregular. Optimization also differs from an uncertainty ensemble: searching for a high apogee does not estimate the spread of apogees under uncertain conditions.

\sourceline{\source{core/src/main/java/info/openrocket/core/optimization/services/DefaultSimulationModifierService.java}{available modifiers}; \source{core/src/main/java/info/openrocket/core/optimization/rocketoptimization/modifiers/GenericComponentModifier.java}{component modifier}; \source{core/src/main/java/info/openrocket/core/optimization/rocketoptimization/domains/StabilityDomain.java}{stability domain}; \source{core/src/main/java/info/openrocket/core/optimization/general/onedim/GoldenSectionSearchOptimizer.java}{golden-section search}; \source{core/src/main/java/info/openrocket/core/optimization/general/multidim/MultidirectionalSearchOptimizer.java}{multidirectional search}.}

